\chapter{Implementation}\label{ch:impl}

This chapter walks through the running system module by module.  Each section
pairs a short code excerpt from the production source tree with the
corresponding piece of the running dashboard, so the reader can trace an idea
from mathematics to Python to a rendered UI.  All listings are verbatim
extracts from files under \texttt{src/}; comments and blank lines were
trimmed only where they carried no explanatory value.

\section{Data Ingestion and Pre-processing}\label{sec:ingest}
Ingestion runs on a scheduler that polls each subreddit's public RSS feed at
a configurable cadence (default: hourly).  Historical dumps are pulled from
the Arctic Shift API for the sentiment gold set.  Each incoming post is
normalised (HTML-stripped, lowercased for lexical filters), deduplicated by
\texttt{post\_id}, and screenshots are downloaded to a local image cache.
Listing~\ref{lst:arctic} shows the paged fetch loop that drives the
historical pull; it uses \texttt{curl} as transport to sidestep TLS
quirks on corporate networks and emits an \texttt{on\_page} progress
callback so the dashboard's Pipeline view can render live counters.
Figure~\ref{fig:ingest} shows the overall flow.

\begin{lstlisting}[caption={Paged Reddit ingestion with cursor-based dedup ({\tt src/ingestion/arctic\_shift.py}).},label={lst:arctic}]
def fetch_posts_arctic(subreddit, since_utc=0.0, limit=500,
                       self_posts_only=False, on_page=None):
    fetched, before = 0, None
    while fetched < limit:
        batch_size = min(100, limit - fetched)
        params = {"subreddit": subreddit,
                  "limit": str(batch_size), "sort": "desc"}
        if since_utc > 0:  params["after"]  = str(int(since_utc))
        if before:         params["before"] = str(int(before))
        if self_posts_only: params["is_self"] = "true"

        data = _curl_get(f"{BASE_URL}?{urlencode(params)}")
        if not data or not (posts := data.get("data", [])):
            break

        for post in posts:
            created = post.get("created_utc", 0)
            if since_utc > 0 and created <= since_utc:
                continue                     # already ingested
            selftext = (post.get("selftext") or "").lower()
            if selftext in ("[deleted]", "[removed]"):
                continue                     # moderator-scrubbed
            if (norm := _normalize_post(post, subreddit)):
                yield norm
                fetched += 1

        if len(posts) < batch_size: break
        before = posts[-1].get("created_utc", 0)  # page cursor
        time.sleep(0.5)                            # be polite
\end{lstlisting}

\begin{figure}[htbp]
\centering\includegraphics[width=\linewidth]{fig9_ingestion_flow}
\caption{Data ingestion and pre-processing flow.}\label{fig:ingest}
\end{figure}

\section{Trust-Score Module}\label{sec:trust}
The trust score $T \in [0,1]$ blends account metadata and content signals:
$T = 0.4\,M + 0.3\,C + 0.3\,L$ where $M$ is a metadata score
(account\_age, karma, verified\_email, cakeday), $C$ is an LLM-graded
credibility signal, and $L$ is a lexical retail-insider score.  Priority
tiers are: P1 iff $T \ge 0.75$ and sentiment is negative with confidence
$\ge 0.75$; P2 iff $T \ge 0.55$ and negative with confidence $\ge 0.65$.
The core scorer is small on purpose --- Listing~\ref{lst:trust} shows the
entire \texttt{TrustScorer.score()} method.  It also implements the
cost-shaping heuristic that only invokes the LLM credibility check when
metadata is genuinely ambiguous ($0.3 < M < 0.8$); outside that band
metadata is already decisive.  Figure~\ref{fig:trust} depicts the
composition.  Every constant is exposed in the dashboard so that a
stakeholder can argue with the rule rather than guess at a black box.

\begin{lstlisting}[caption={Combined trust scorer with cost-shaped LLM invocation ({\tt src/trust/scorer.py}).},label={lst:trust}]
def score(self, unit: dict) -> dict:
    meta  = score_metadata(unit, self.config)     # M
    dedup = score_originality(unit)               # L

    # LLM credibility (C) --- only when metadata is ambiguous.
    # Outside the band, metadata is decisive, so we skip the call
    # to cap cost on cloud providers.
    llm_score, llm_flags = 0.5, []
    if self.llm_client and 0.3 < meta < 0.8:
        r = self.llm_client.check_credibility(
            text=self._get_text(unit),
            metadata=unit.get("author_metadata", {}))
        llm_score = r.get("credibility_score", 0.5)
        llm_flags = r.get("flags", [])

    # T = 0.4 M + 0.3 dedup + 0.3 LLM, clamped to [0, 1]
    combined = (self.config.weight_metadata * meta   +
                self.config.weight_dedup    * dedup  +
                self.config.weight_llm      * llm_score)
    combined = round(min(max(combined, 0.0), 1.0), 3)

    unit["trust_score"]      = combined
    unit["trust_components"] = {"metadata": meta,
                                "dedup":    dedup,
                                "llm":      llm_score}
    unit["trust_flags"]      = llm_flags
    unit["is_trusted"]       = combined >= self.config.threshold
    if not unit["is_trusted"]:
        unit["processing_status"] = "flagged_low_trust"
    return unit
\end{lstlisting}

\begin{figure}[htbp]
\centering\includegraphics[width=\linewidth]{fig3_trust_composition}
\caption{Trust-Score composition and P1/P2 tier rules.}\label{fig:trust}
\end{figure}

\section{Sentiment Classification (ModernBERT)}\label{sec:sent}
The sentiment head is a fine-tuned ModernBERT-base with a three-class
softmax (positive, neutral, negative).  Training used a three-stage
curriculum: (Stage 1) 8k pseudo-labelled tweets from the twitter-roberta
teacher, (Stage 2) 2k retail-focused Reddit posts pseudo-labelled by the
same teacher, and (Stage 3) 200 hand-labelled retail Reddit posts for a
final domain adaptation.  See Figure~\ref{fig:curriculum}.  The 5-fold
out-of-fold macro-F1 rose from the baseline of 0.6272 to \textbf{0.7642};
long-post accuracy (bucket $\ge 512$~tokens, $n{=}7$, all negative-class)
rose from 5/7 correct on the RoBERTa baseline to \textbf{7/7 correct}
on ModernBERT.

At inference time the sentiment head is wrapped by a thin
\texttt{SentimentAnalyzer} that owns exactly two responsibilities: assemble
the prompt string (subreddit-aware, and injecting an
\texttt{[image: \dots]} tail when a Gemma caption is present), and shape
the model's output into the analysis record the storage layer expects.
Listing~\ref{lst:analyzer} shows both.

\begin{lstlisting}[caption={Sentiment orchestration with vision-caption fusion ({\tt src/analysis/analyzer.py}).},label={lst:analyzer}]
def _get_text(self, unit: dict) -> str:
    """Compose a subreddit-tagged prompt; append the Gemma caption
    as [image: ...] so downstream models see it as post body."""
    title = unit.get("title", "") or ""
    body  = unit.get("body",  "") or ""
    sub   = unit.get("subreddit", "unknown")
    cap   = (unit.get("image_caption") or "").strip()
    cap_suffix = f" [image: {cap}]" if cap else ""

    if unit.get("unit_type") == "comment":
        return f'Comment from r/{sub}: "{body}"'
    if title and body:
        return f'Post from r/{sub}: "{title}" -- {body}{cap_suffix}'
    if title:
        return f'Post from r/{sub}: "{title}"{cap_suffix}'
    if body:
        return f'Post from r/{sub}: "{body}{cap_suffix}"'
    # Image-only post: the caption IS the content.
    return f'Post from r/{sub}: "{cap}"' if cap else f'Post from r/{sub}: ""'

def _build_analysis_record(self, unit, result):
    unit_id = unit.get("id", "")
    conf    = result.get("sentiment_confidence", 0.0)
    return {
        "id": f"analysis_{unit_id}",  "post_id": unit_id,
        "subreddit":            unit.get("subreddit", ""),
        "trust_score":          unit.get("trust_score"),
        "sentiment":            result.get("sentiment", "neutral"),
        "sentiment_confidence": conf,
        "aspects":              result.get("aspects", []),
        "key_phrases":          result.get("key_phrases", []),
        "summary":              result.get("summary", ""),
        "model_used":     result.get("model_used", self.llm.model_name),
        "model_version":  result.get("model_version", self.llm.model_version),
        "analyzed_at":    datetime.now(timezone.utc).isoformat(),
        "needs_review":   conf < self.config.confidence_threshold,
        "partition_key":  unit.get("subreddit", ""),
    }
\end{lstlisting}

\begin{figure}[htbp]
\centering\includegraphics[width=\linewidth]{fig10_modernbert_curriculum}
\caption{ModernBERT three-stage fine-tuning curriculum.}\label{fig:curriculum}
\end{figure}

\section{Aspect-Based Opinion Mining}\label{sec:aspect}
Aspects are extracted with a zero-shot NLI classifier
(\texttt{mDeBERTa-v3-base-xnli-multilingual-nli-2mil7}) scored against a fixed
eight-aspect taxonomy (Table~\ref{tab:aspects}).  A post can carry multiple
aspects when their entailment scores are above threshold~0.55.
Figure~\ref{fig:aspect} illustrates the NLI framing.  On cloud LLM
providers the same taxonomy is enforced with the structured JSON prompt in
Listing~\ref{lst:aspect-prompt}: the model is told exactly which aspect
group applies to which subreddit (customer vs.\ employee), which prevents
the ``rude cashier'' $\leftrightarrow$ ``bad manager'' bleed that
plagued the first-cut prompts.

\begin{lstlisting}[caption={Aspect taxonomy enforced in the JSON-mode system prompt ({\tt src/analysis/prompts.py}).},label={lst:aspect-prompt}]
SENTIMENT_ASPECT_SYSTEM_PROMPT = """You are a retail analyst ...

Rules:
1. Sentiment must be one of: "positive", "negative", "neutral".
2. Choose aspects from the correct group based on the subreddit:

   CUSTOMER aspects (r/walmart, r/samsclub, r/WalmartPlus, ...):
     store_experience, online_app, delivery_pickup, product_quality,
     returns, customer_support, pricing

   EMPLOYEE aspects (r/WalmartEmployees, r/OGPBackroom,
                     r/Sparkdriver, r/walmart_RX, ...):
     workforce_hr, pay_benefits, management, safety_policy, workload

3. A post can have 0 aspects (off-topic) or multiple aspects.
4. NEVER mix customer and employee aspects in the same post.
5. Handle sarcasm and slang (Reddit tone). "Love how my order
   arrived 3 days late" = negative + delivery_pickup.
6. Confidence scores (0.0-1.0) for each prediction.

Output JSON:
{
  "sentiment": "positive|negative|neutral",
  "sentiment_confidence": 0.0-1.0,
  "aspects": [
    {"aspect": "<name>", "sentiment": "...", "confidence": 0.0-1.0}
  ],
  "key_phrases": ["phrase1", "phrase2"],
  "summary": "One sentence summary of the issue/praise"
}"""
\end{lstlisting}

\begin{table}[htbp]\centering
\caption{Eight-aspect retail taxonomy.}\label{tab:aspects}
\small
\begin{tabular}{@{}ll@{}}\toprule
\bfseries Aspect & \bfseries Example indicative terms\\\midrule
Pricing               & ``\$2 more this week'', ``price hike'', ``clearance''\\
Product quality       & ``broken'', ``expired'', ``mouldy''\\
Customer service      & ``rude cashier'', ``no help'', ``manager''\\
Store experience      & ``dirty floor'', ``long line'', ``restroom''\\
Online / app          & ``app crash'', ``payment failed'', ``2FA loop''\\
Delivery / pickup     & ``missed slot'', ``OGP'', ``Spark driver''\\
Returns               & ``refused return'', ``restocking fee''\\
Account / login       & ``locked out'', ``password reset''\\\bottomrule
\end{tabular}
\end{table}

\begin{figure}[htbp]
\centering\includegraphics[width=0.95\linewidth]{fig11_aspect_nli}
\caption{Zero-shot NLI aspect tagging.}\label{fig:aspect}
\end{figure}

\section{Vision / Image Processing}\label{sec:vision}
Screenshots are captioned by Gemma~3 4B via Ollama using a multi-pass prompt
pipeline: Pass~1 identifies the image \emph{structure}
(screenshot / receipt / product page / error dialog / meme / photo);
Pass~2 tiles the image at native resolution and extracts text per tile;
Pass~3 merges the tile texts against the structure into 2--4 factual
sentences with a strict ``do not invent'' clause.  A circuit breaker
disables the whole vision branch after three consecutive Ollama failures
so a stopped daemon cannot flood the logs.
Figure~\ref{fig:vision} shows the pipeline.  Hallucination rate on 32
retail screenshots dropped from 50\% (single pass) to \textbf{0\%}; text
extraction rose from 25\% to \textbf{75\%}.
Listing~\ref{lst:vision} shows the entry point and the merge prompt that
carries the anti-hallucination clause.

\begin{lstlisting}[caption={Vision multi-pass caption entry point and merge prompt ({\tt src/analysis/vision.py}).},label={lst:vision}]
MERGE_PROMPT = """Combine these observations about a Walmart-related
image into 2-4 clear sentences.

Image type: {structure}
Text found in different regions:
{tile_texts}

Write a factual description that captures: what type of
screen/image this is, all important text (especially errors,
prices, status messages), and what the customer might be
complaining about. Do NOT invent details not mentioned above."""

def caption_enhanced(self, image_path: Path) -> str:
    if not self.cfg.enabled or not Path(image_path).exists():
        return ""
    img = Image.open(image_path)

    # Pass 1: structure identification
    full_b64  = self._image_to_b64(img)
    structure = self._call_ollama(self.model, STRUCTURE_PROMPT, full_b64)
    if not structure:
        return self.caption(image_path)           # fall back

    # Only tile screenshots / receipts / app screens; photos and
    # memes are fine as single-pass.
    if not any(kw in structure.lower() for kw in
               ("screenshot","receipt","app_screen",
                "error_dialog","document","product_page")):
        return self.caption(image_path)

    # Pass 2 + 3: tile at native resolution and extract text per tile
    tiles, tile_texts = self._create_tiles(img), []
    for i, b64 in enumerate(tiles):
        t = self._call_ollama(self.model, TILE_TEXT_PROMPT, b64)
        if t and "NO TEXT" not in t.upper():
            tile_texts.append(f"Region {i+1}: {t}")

    merged = self._call_ollama(
        self.model,
        MERGE_PROMPT.format(structure=structure,
                            tile_texts="\n".join(tile_texts)),
        full_b64)
    return merged or self.caption(image_path)
\end{lstlisting}

\begin{figure}[htbp]
\centering\includegraphics[width=\linewidth]{fig12_vision_multipass}
\caption{Vision multi-pass anti-hallucination pipeline.}\label{fig:vision}
\end{figure}

\section{Aggregation, Alerts and Notifications}\label{sec:agg}
The aggregator writes daily brand-health rollups (sentiment share, top
aspects, trust bucket) into SQLite.  The alert engine fires on volume,
sentiment-crash, emerging-topic and competitor rules; each rule reads its
thresholds from \texttt{data/alert\_rules.json}, which the dashboard's
Alert Rules panel writes to.  Listing~\ref{lst:alerts} shows the two
detectors that drive the majority of P1/P2 traffic: a volume-spike detector
that fires when today's post count crosses $\mu + \sigma\!\cdot\!\theta$ of
the trailing 7-day mean, and a sentiment-crash detector that fires when
the negative share jumps by more than 0.30 versus yesterday.  Alerts flow
to the dashboard's Alert Feed (\S\ref{sec:dash}) and, optionally, to an
email digest.  Figure~\ref{fig:alerts} shows the routing.

\begin{lstlisting}[caption={Volume-spike and sentiment-crash detectors ({\tt src/alerts/engine.py}).},label={lst:alerts}]
def detect_volume_spike(self, sigma_threshold: float = 2.0):
    """Fire when today's post count is > mu + sigma * threshold of
       the trailing 7-day mean."""
    now   = datetime.now(timezone.utc)
    today = now.strftime("%Y-%m-%d")
    today_agg = self.storage.get_item("aggregates",
                                      f"agg_{today}_daily", today)
    if not today_agg: return []

    counts = []
    for i in range(1, 8):
        d = (now - timedelta(days=i)).strftime("%Y-%m-%d")
        past = self.storage.get_item("aggregates", f"agg_{d}_daily", d)
        if past: counts.append(past.get("total_posts", 0))
    if len(counts) < 3: return []

    mu  = sum(counts) / len(counts)
    std = math.sqrt(sum((x-mu)**2 for x in counts) / len(counts))
    if std == 0: return []

    n = today_agg.get("total_posts", 0)
    z = (n - mu) / std
    if z <= sigma_threshold: return []
    return [{
        "id": f"alert_spike_{today}", "type": "volume_spike",
        "severity": "high" if z > 3.0 else "medium",
        "title": f"Volume spike: {n} posts today "
                 f"({z:.1f}sigma above mean)",
        "details": {"today_count": n, "mean_7d": round(mu, 1),
                    "std_7d":  round(std, 1),
                    "z_score": round(z, 2)},
        "detected_at": now.isoformat(), "time_window": today,
    }]

def detect_sentiment_crash(self, drop_threshold: float = 0.3):
    """Fire when negative share jumps > 0.3 vs yesterday."""
    now   = datetime.now(timezone.utc)
    today = now.strftime("%Y-%m-%d")
    yday  = (now - timedelta(days=1)).strftime("%Y-%m-%d")
    t = self.storage.get_item("aggregates", f"agg_{today}_daily", today)
    y = self.storage.get_item("aggregates", f"agg_{yday}_daily",  yday)
    if not t or not y: return []
    delta = (t.get("sentiment_distribution", {}).get("negative", 0)
           - y.get("sentiment_distribution", {}).get("negative", 0))
    if delta <= drop_threshold: return []
    return [{
        "id": f"alert_crash_{today}", "type": "sentiment_crash",
        "severity": "critical" if delta > 0.4 else "high",
        "title": f"Sentiment crash: negative +{delta:.0%} vs yesterday",
        "details": {"delta": round(delta, 3)},
        "detected_at": now.isoformat(), "time_window": today,
    }]
\end{lstlisting}

\begin{figure}[htbp]
\centering\includegraphics[width=\linewidth]{fig13_alert_routing}
\caption{Alert engine and notification routing.}\label{fig:alerts}
\end{figure}

The rendered Alert Feed is shown in Figure~\ref{fig:ui:alerts}; each row
is one entry from the JSON returned by the detectors above.

\begin{figure}[htbp]
\centering\includegraphics[width=\linewidth,height=0.82\textheight,keepaspectratio]{ui/alert_feed.png}
\caption{Alert Feed --- live P1/P2 stream sourced from the volume-spike
and sentiment-crash detectors of Listing~\ref{lst:alerts}.}\label{fig:ui:alerts}
\end{figure}

\section{Dashboard Overview}\label{sec:dash}
The dashboard is a Vite + React 18 + Tailwind SPA served by FastAPI.
Pages: Brand Health (KPIs and sentiment trend), Alert Feed (live P1/P2
stream), Post Explorer (multi-facet search), Review \& Validate (HITL
queue), Post Lifecycle (Kanban), Insights \& Competitor, Pipeline
Control, and Notifications.  The landing page is Brand Health
(Figure~\ref{fig:ui:brand}); the entire volume trend, sentiment gauge
and aspect breakdown are driven by the SQLite aggregates that
Listing~\ref{lst:alerts} also reads from, so what the analyst sees on
the dashboard is always the same view of the world that the alert engine
scored.

\begin{figure}[htbp]
\centering\includegraphics[width=\linewidth,height=0.82\textheight,keepaspectratio]{ui/brand_health.png}
\caption{Brand Health --- sentiment gauge, distribution donut, volume trend
and priority breakdown driven by the daily aggregates.}\label{fig:ui:brand}
\end{figure}

\section{Human-in-the-Loop Review \& Validate}\label{sec:hitl}
Every post landing in the review queue can be re-labelled (sentiment,
aspects), corrected (trust bucket), replied to (rule or LLM draft), and
closed with a resolution reason.  The workflow is shown in
Figure~\ref{fig:hitl}; the rendered queue is in
Figure~\ref{fig:ui:review}.  All actions are logged as
(post, before, after, analyst\_id, timestamp).  Listing~\ref{lst:review}
shows the endpoint that receives an analyst correction: it appends an
immutable feedback record \emph{and} rewrites the analysis so the
correction flows through to every downstream view --- Brand Health,
aspect drilldowns, Post Explorer and the Kanban board --- without any
extra re-aggregation call.

\begin{lstlisting}[caption={HITL correction endpoint: feedback audit log + live analysis rewrite ({\tt src/dashboard/api.py}).},label={lst:review}]
@app.post("/api/review/{post_id}")
def submit_review(post_id: str, correction: dict):
    _ensure_initialized()
    now = datetime.now(timezone.utc).isoformat()

    # 1. Immutable audit log (feedback table).
    feedback = {
        "id": f"fb_{post_id}_"
              f"{int(datetime.now(timezone.utc).timestamp())}",
        "post_id":             post_id,
        "analyst_id":          correction.get("analyst_id", "default"),
        "original_sentiment":  correction.get("original_sentiment"),
        "corrected_sentiment": correction.get("corrected_sentiment"),
        "original_aspects":    correction.get("original_aspects", []),
        "corrected_aspects":   correction.get("corrected_aspects", []),
        "trust_override":      correction.get("trust_override"),
        "notes":               correction.get("notes", ""),
        "created_at":          now,
        "partition_key":       correction.get("analyst_id", "default"),
    }
    _storage.upsert("feedback", feedback)

    # 2. Rewrite the analysis so the correction flows through to
    #    Brand Health, aspect drilldowns and the Kanban board.
    analysis = _storage.get_item("analyses",
                 f"analysis_{post_id}", correction.get("subreddit", ""))
    if analysis:
        if (c := correction.get("corrected_sentiment")):
            analysis["sentiment"]            = c
            analysis["sentiment_confidence"] = 1.0     # human-validated
        if isinstance(correction.get("corrected_aspects"), list):
            analysis["aspects"] = [
                {"aspect": a, "confidence": 1.0} if isinstance(a, str)
                else a for a in correction["corrected_aspects"]]
        if (t := correction.get("trust_override")) is not None:
            analysis["trust_score"]    = max(0.0, min(1.0, float(t)))
            analysis["trust_override"] = True
        analysis["needs_review"]    = False
        analysis["human_validated"] = True
        analysis["validated_at"]    = now
        analysis["validated_by"]    = correction.get("analyst_id", "default")
        _storage.upsert("analyses", analysis)

    return {"status": "saved", "feedback_id": feedback["id"],
            "analysis_updated": analysis is not None}
\end{lstlisting}

\begin{figure}[htbp]
\centering\includegraphics[width=\linewidth]{fig5_review_validate_flow}
\caption{Human-in-the-loop Review \& Validate flow.}\label{fig:hitl}
\end{figure}

\begin{figure}[htbp]
\centering\includegraphics[width=\linewidth,height=0.82\textheight,keepaspectratio]{ui/review_validate.png}
\caption{Review \& Validate queue --- posts flagged by the model land here;
analysts re-label sentiment/aspects and either compose a reply or resolve
in place.  Corrections are POSTed to the endpoint in
Listing~\ref{lst:review}.}\label{fig:ui:review}
\end{figure}

\section{Smart Reply Composer --- Dual-Draft (Rule + LLM)}\label{sec:reply}
The reply composer offers two drafts side-by-side: a rule-based template
that guarantees policy compliance, and an LLM-generated draft that
personalises to the post's aspects and sentiment.  Few-shot examples come
from the last analyst-posted replies pulled straight from the feedback
store --- this is what closes the learning loop of \S\ref{sec:loop}
without any offline retraining.  Listing~\ref{lst:reply} shows the endpoint
that ties it together.  See Figure~\ref{fig:reply}.

\begin{lstlisting}[caption={Dual-draft reply endpoint (rule + LLM, few-shot from posted replies) ({\tt src/dashboard/api.py}).},label={lst:reply}]
@app.post("/api/review/{post_id}/draft-reply")
def draft_reply(post_id: str, payload: dict | None = None):
    """Generate a customer-specific reply draft. Uses the last N
    analyst-posted replies as few-shot examples --- every posted
    reply becomes training context for future drafts."""
    _ensure_initialized()
    payload = payload or {}
    subredd = payload.get("subreddit", "")

    analysis = _storage.get_item("analyses",
                                 f"analysis_{post_id}", subredd)
    if not analysis:
        return {"status": "error", "reason": "analysis_not_found"}
    if analysis.get("sentiment") != "negative":
        return {"status": "error", "reason": "not_negative"}

    raw      = _storage.get_item("raw_posts", post_id, subredd) or {}
    aspects  = _aspect_names(analysis.get("aspects", []))
    examples = _collect_reply_examples(limit=5)      # <-- few-shot

    llm    = _get_reply_llm()
    result = llm.generate_reply_pair(
        post_title=raw.get("title", ""),
        post_text =raw.get("body",  "") or raw.get("title", ""),
        subreddit =analysis.get("subreddit", ""),
        author    =raw.get("author", ""),
        aspects   =aspects,
        examples  =examples)

    drafts  = result.get("drafts", [])       # rule + LLM
    primary = drafts[0] if drafts else {}
    return {"status": "ok", "drafts": drafts,
            "reply":         primary.get("reply", ""),
            "model_used":    primary.get("model_used", ""),
            "source":        primary.get("source", ""),
            "examples_used": len(examples)}
\end{lstlisting}

\begin{figure}[htbp]
\centering\includegraphics[width=\linewidth]{fig6_reply_cascade}
\caption{Smart Reply Composer --- dual-draft with few-shot.}\label{fig:reply}
\end{figure}

\section{Learning Loop --- Feedback for Future Retraining}\label{sec:loop}
Corrections and posted replies stream into the feedback store keyed by
post\_id and model version (the same audit log written by
Listing~\ref{lst:review}).  A weekly export produces training pairs for
sentiment (correction $\to$ label) and reply generation (post $\to$
posted-reply).  Figure~\ref{fig:loop} depicts the loop and the path to
future auto-reply --- the same few-shot pull in Listing~\ref{lst:reply}
means the composer becomes measurably better every week without touching
the model weights.

\begin{figure}[htbp]
\centering\includegraphics[width=\linewidth]{fig7_learning_loop}
\caption{Feedback learning loop and path to auto-reply.}\label{fig:loop}
\end{figure}

\section{Post Explorer and Multi-Facet Search}\label{sec:explore}
Post Explorer supports filters over subreddit, aspect, sentiment, trust
bucket, priority tier, resolution state, date range, and free-text.
Figure~\ref{fig:ui:posts} shows the rendered page; every filter is a
query parameter on the same \texttt{/api/posts} endpoint that the Kanban
and Review queues also read, so an analyst can navigate freely between
the three views without ever seeing an inconsistent post state.

\begin{figure}[htbp]
\centering\includegraphics[width=\linewidth,height=0.82\textheight,keepaspectratio]{ui/post_explorer.png}
\caption{Post Explorer --- multi-facet search over subreddit, aspect,
sentiment, trust bucket, priority tier and free text.}\label{fig:ui:posts}
\end{figure}

\section{Post Lifecycle (Kanban Workflow)}\label{sec:kanban}
Every P1/P2 alert flows through a Kanban board with four states
(Triaged $\to$ Acknowledged $\to$ In Progress $\to$ Resolved).  See
Figure~\ref{fig:kanban} for the state machine and
Figure~\ref{fig:ui:kanban} for the rendered board.  State transitions
carry SLA timestamps so we can compute time-to-acknowledge and
time-to-resolve.

\begin{figure}[htbp]
\centering\includegraphics[width=\linewidth]{fig8_lifecycle_kanban}
\caption{Post Lifecycle Kanban workflow.}\label{fig:kanban}
\end{figure}

\begin{figure}[htbp]
\centering\includegraphics[width=\linewidth,height=0.82\textheight,keepaspectratio]{ui/lifecycle_kanban.png}
\caption{Post Lifecycle --- Kanban board with SLA timestamps for every
transition.}\label{fig:ui:kanban}
\end{figure}

\section{Insights and Competitor Analysis}\label{sec:insights}
The Insights page shows aspect-level negative-share for Walmart vs.\ a
curated set of competitor communities (Target, Costco, Amazon, Kroger,
Aldi) so analysts can distinguish a company-specific issue from an
industry trend.  Figure~\ref{fig:ui:insights} shows the rendered view.

\begin{figure}[htbp]
\centering\includegraphics[width=\linewidth,height=0.82\textheight,keepaspectratio]{ui/insights_competitor.png}
\caption{Insights \& Competitor --- aspect-level negative-share for Walmart
vs.\ curated competitor communities.}\label{fig:ui:insights}
\end{figure}

Notifications route the same alerts to email digests and the in-app
notification centre.  The rendered notification centre is shown in
Figure~\ref{fig:ui:notif}.

\begin{figure}[htbp]
\centering\includegraphics[width=\linewidth,height=0.82\textheight,keepaspectratio]{ui/notifications.png}
\caption{Notification centre --- in-app mirror of the email digest;
sourced from the same alert rows that drive
Figure~\ref{fig:ui:alerts}.}\label{fig:ui:notif}
\end{figure}

\section{Storage Layer}\label{sec:storage}
Persistence uses SQLite in WAL mode for hot data (posts, aggregates,
alerts, feedback), plus a JSONL cost log for LLM calls and a local file
cache for downloaded images.  A backup runs nightly.  The schema is
declared in a single \texttt{CREATE TABLE} block so a fresh clone of the
repository can boot with an empty database in one command;
Listing~\ref{lst:storage} shows the relevant excerpt --- note the
\texttt{data JSON} column that holds the full document, mirroring the
Cosmos DB partition/PK design so we can lift-and-shift to Azure
without changing calling code.

\begin{lstlisting}[caption={SQLite schema mirroring the Cosmos DB partition design ({\tt src/storage/store.py}).},label={lst:storage}]
class SQLiteBackend(StorageBackend):
    def __init__(self, db_path="data/local.db"):
        self.db_path = Path(db_path)
        self.db_path.parent.mkdir(parents=True, exist_ok=True)
        self._conn = sqlite3.connect(str(self.db_path),
                                     check_same_thread=False)
        self._conn.row_factory = sqlite3.Row
        self._init_tables()

    def _init_tables(self):
        self._conn.executescript("""
            CREATE TABLE IF NOT EXISTS raw_posts (
                id TEXT PRIMARY KEY, subreddit TEXT,
                data JSON NOT NULL, created_timestamp REAL,
                processing_status TEXT DEFAULT 'pending');
            CREATE INDEX IF NOT EXISTS idx_raw_posts_subreddit
                ON raw_posts(subreddit);
            CREATE INDEX IF NOT EXISTS idx_raw_posts_status
                ON raw_posts(processing_status);

            CREATE TABLE IF NOT EXISTS analyses (
                id TEXT PRIMARY KEY, post_id TEXT,
                subreddit TEXT, data JSON NOT NULL);
            CREATE INDEX IF NOT EXISTS idx_analyses_post
                ON analyses(post_id);

            CREATE TABLE IF NOT EXISTS aggregates (
                id TEXT PRIMARY KEY, time_window TEXT,
                window_type TEXT, data JSON NOT NULL);

            CREATE TABLE IF NOT EXISTS feedback (
                id TEXT PRIMARY KEY, analyst_id TEXT,
                post_id TEXT, data JSON NOT NULL);

            CREATE TABLE IF NOT EXISTS alerts (
                id TEXT PRIMARY KEY, type TEXT, severity TEXT,
                time_window TEXT, data JSON NOT NULL);
        """)
        self._conn.execute("PRAGMA journal_mode=WAL;")
        self._conn.commit()
\end{lstlisting}
